% SGA 3, Expose VIII, English reconstruction.
% Scope: Sections 2 and 3 in full, and Section 4 from its heading through
% the word "Moreover," in the proof of Corollary 4.7.
% Authority: Polo--Gille Exp8-8nov09.pdf, standalone pages 7--13
% (combined-reader pages 623--629).
% Jacob Reinhold's CC BY 4.0 Markdown lineage was used only as a phrasing
% comparison. The Polo--Gille PDF controls prose, formulas, numbering,
% notes, references, and diagrams.
%
% Start cursor: "2. Proprietes schematiques des groupes diagonalisables"
% and "Elles sont resumees dans la" (standalone p. 7).
% End cursor: "De plus," in the proof of Corollary 4.7
% (standalone p. 13); the sentence continues on standalone p. 14.

\setcounter{footnote}{10}

% Source page: Exp8 standalone p. 7; combined-reader p. 623.
\SGARefTarget{sga3:viii:section:2}\section*{2. 可对角化群的概形论性质}
\addcontentsline{toc}{section}{2. 可对角化群的概形论性质}

这些性质概括于\SGARefLink{sga3:viii:proposition:2:1}{下述命题}。

% Source page: Exp8 standalone p. 8; combined-reader p. 624.
\medskip
\noindent\SGARefTarget{sga3:viii:proposition:2:1}\textbf{命题 2.1.}
\textit{设 \(S\) 为非空预概形，\(M\) 为通常的交换群，\(G=D(M_S)\) 为 \(M\) 定义的可对角化 \(S\)-群。则：}
\begin{enumerate}
  \item[(a)]\SGARefTarget{sga3:viii:proposition:2:1:item:a:01}
  \textit{\(G\) 在 \(S\) 上忠实平坦，且在 \(S\) 上仿射（从而更是在 \(S\) 上拟紧）；}

  \item[(b)]\SGARefTarget{sga3:viii:proposition:2:1:item:b:01}
  \textit{\(M\) 为有限型，当且仅当 \(G\) 在 \(S\) 上有限型，也当且仅当 \(G\) 在 \(S\) 上有限呈示；}

  \item[(c)]\SGARefTarget{sga3:viii:proposition:2:1:item:c:01}
  \textit{\(M\) 有限，当且仅当 \(G\) 在 \(S\) 上有限，也当且仅当 \(G\) 在 \(S\) 上有限型且被某个整数 \(n>0\) 零化。此时，}
  \[
    \deg(G/S)=\operatorname{Card}(M);
  \]

  \item[\((c')\)]\SGARefTarget{sga3:viii:proposition:2:1:item:cprime:01}
  \textit{\(M\) 为挠群，当且仅当 \(G\) 在 \(S\) 上整；}

  \item[(d)]\SGARefTarget{sga3:viii:proposition:2:1:item:d:01}
  \textit{\(M=0\) 当且仅当 \(G\) 为单位 \(S\)-群；}

  \item[(e)]\SGARefTarget{sga3:viii:proposition:2:1:item:e:01}
  \textit{\(M\) 为有限型且其挠子群的阶与 \(S\) 的各剩余特征互素，当且仅当 \(G\) 在 \(S\) 上平滑。}
\end{enumerate}

\SGARefLink{sga3:viii:proposition:2:1:item:a:01}{(a)} 至 \SGARefLink{sga3:viii:proposition:2:1:item:d:01}{(d)} 的验证直接，留给读者。来证明 \SGARefLink{sga3:viii:proposition:2:1:item:e:01}{(e)}。若 \(G\) 在 \(S\) 上平滑，则它在 \(S\) 上局部有限呈示；由于它在 \(S\) 上仿射，故在 \(S\) 上有限呈示，因此 \(M\) 为有限型。所以一开始就可假设 \(M\) 为有限型，从而 \(G\) 在 \(S\) 上有限呈示。此时\footnote{编者注：这里用到了由 \SGARefLink{sga3:viii:proposition:2:1:item:a:01}{(a)} 得到的 \(G\) 在 \(S\) 上平坦这一事实。}，\(G\) 在 \(S\) 上平滑当且仅当其几何纤维平滑，这就归结到 \(S\) 为代数闭域 \(k\) 的谱的情形。

写成 \(M=T\times L\)，其中 \(T\) 为挠子群，\(L\) 为自由群，且 \(L\simeq\mathbb Z^r\)。则
\[
  D(M)=D(T)\times D(L),
  \qquad
  D(L)\simeq\mathbb G_m^r,
\]
且 \(D(L)\) 在 \(k\) 上平滑。因此，\(G=D(M)\) 在 \(k\) 上平滑当且仅当 \(D(T)\) 平滑。由于 \(D(T)\) 在 \(k\) 上有限，其次数等于 \(T\) 的阶 \(n\)，这意味着 \(D(T)(k)\) 有 \(n\) 个元素。而 \(T\) 同构于若干群 \(\mathbb Z/n_i\mathbb Z\) 的和，其中 \(n=\prod_i n_i\)；从而
\[
  D(T)=\prod_iD(\mathbb Z/n_i\mathbb Z)
      =\prod_i\mu_{n_i},
\]
其中 \(\mu_{n_i}\) 为 \(n_i\) 次单位根群概形。因此
\[
  \operatorname{card}\bigl(D(T)(k)\bigr)
  =\prod_i\operatorname{card}\bigl(\mu_{n_i}(k)\bigr).
\]
这里 \(\operatorname{card}(\mu_{n_i}(k))\) 为 \(k\) 中 \(n_i\) 次单位根的数目，至多为 \(n_i\)，且等号成立当且仅当 \(n_i\) 与 \(k\) 的特征 \(p\) 互素。
\emph{〔原文两处都印作 \(\mu_i(k)\)。〕}
由此可知，
\[
  \operatorname{card}\bigl(D(T)(k)\bigr)=n,
  \qquad n=\prod_i n_i,
\]
当且仅当每个 \(n_i\) 都与 \(p\) 互素，即当且仅当 \(n\) 与 \(p\) 互素。
\hfill\(\square\)

\SGARefTarget{sga3:viii:section:3}\section*{3. 函子 \texorpdfstring{\(D_S\)}{D\_S} 的正合性质}
\addcontentsline{toc}{section}{3. 函子 \texorpdfstring{\(D_S\)}{D\_S} 的正合性质}

\medskip
\noindent\SGARefTarget{sga3:viii:theorem:3:1}\textbf{定理 3.1.}
\textit{设 \(S\) 为预概形，}
\SGARefTarget{sga3:viii:theorem:3:1:diagram:01}
\[
  0\longrightarrow M'
   \xrightarrow{u}M
   \xrightarrow{v}M''
   \longrightarrow0
\]
\textit{为通常的交换群的正合列。考虑转置同态列}
\SGARefTarget{sga3:viii:theorem:3:1:diagram:02}
\[
  0\longrightarrow D_S(M'')
   \xrightarrow{v^t}D_S(M)
   \xrightarrow{u^t}D_S(M')
   \longrightarrow0.
\]
\begin{enumerate}
  \item[(i)]\SGARefTarget{sga3:viii:theorem:3:1:item:i:01}
  \textit{同态 \(v^t\) 诱导 \(D_S(M'')\) 到 \(u^t\) 的核的同构，且 \(u^t\) 忠实平坦并拟紧。}

% Source page: Exp8 standalone p. 9; combined-reader p. 625.
  \item[(ii)]\SGARefTarget{sga3:viii:theorem:3:1:item:ii:01}%
  \footnote{编者注：增加了以下内容，并相应扩充了证明。}
  \textit{\(D_S(M')\) 表示关于 \((\mathrm{fpqc})\) 拓扑的商层 \(D_S(M)/D_S(M'')\)。}
\end{enumerate}

令 \(\mathcal M\) 表示忠实平坦拟紧态射族。首先，\SGARefLink{sga3:viii:theorem:3:1:item:ii:01}{(ii)} 可由 \SGARefLink{sga3:viii:theorem:3:1:item:i:01}{(i)} 推出（参见 IV，\SGARefLink{sga3:iv:corollary:4:6:5:1}{4.6.5.1}）。事实上，\(u^t\) 在 \(D_S(M)\) 上定义的等价关系，与子群
\[
  \operatorname{Ker}(u^t)=D_S(M'').
\]
定义的等价关系相同。由于 \(u^t\in\mathcal M\)，此等价关系为 \(\mathcal M\)-有效的（参见 IV，\SGARefLink{sga3:iv:scholium:3:3:2:1}{3.3.2.1}）；因此 \(D_S(M')\) 表示关于 \((\mathrm{fpqc})\) 拓扑的商层（参见 IV，\SGARefLink{sga3:iv:lemma:4:6:5}{4.6.5}）。

\SGARefLink{sga3:viii:theorem:3:1:item:i:01}{(i)} 的第一个断言直接来自函子 \(D_S(-)\) 的定义。更一般地，每个左端不带零的正合列
\SGARefTarget{sga3:viii:theorem:3:1:diagram:03}
\[
  M'\longrightarrow M\longrightarrow M''\longrightarrow0
\]
都给出转置正合列
\SGARefTarget{sga3:viii:theorem:3:1:diagram:04}
\[
  0\longrightarrow D_S(M'')
   \longrightarrow D_S(M)
   \longrightarrow D_S(M').
\]
（这在 \(\SGARefLink{sga3:viii:section:1}{\S1}\) 开头所考虑的更一般背景下仍成立。）另一方面，由于 \(D_S(M)\) 和 \(D_S(M')\) 都在 \(S\) 上仿射，不论同态 \(u:M'\to M\) 如何，态射 \(u^t\) 必为仿射，从而更是拟紧。因此，\SGARefLink{sga3:viii:theorem:3:1:item:i:01}{(i)} 的第二个断言可由\SGARefLink{sga3:viii:corollary:3:2}{下面推论}的\SGARefLink{sga3:viii:corollary:3:2:item:a:01}{(a) 项}得出。

\medskip
\noindent\SGARefTarget{sga3:viii:corollary:3:2}\textbf{推论 3.2.}
\textit{设 \(S\) 为非空预概形，\(u:M'\to M\) 为通常的交换群的同态，\(u^t:G\to G'\) 为其转置同态。则：}
\begin{enumerate}
  \item[(a)]\SGARefTarget{sga3:viii:corollary:3:2:item:a:01}
  \textit{\(u\) 为单态射，当且仅当 \(u^t\) 忠实平坦；}

  \item[(b)]\SGARefTarget{sga3:viii:corollary:3:2:item:b:01}
  \textit{\(u\) 为满态射，当且仅当 \(u^t\) 为单态射；此时 \(u^t\) 实际为闭浸入。}
\end{enumerate}

为证明 \SGARefLink{sga3:viii:corollary:3:2:item:a:01}{(a)}，注意若 \(u\) 为单态射，则 \(\mathcal O_S(M)\) 作为 \(\mathcal O_S(M')\) 上的模有一组非空基：即 \(M\) 模 \(M'\) 的任一代表系所定义的截面系。因此它忠实平坦。反过来，若 \(\mathcal O_S(M)\) 在 \(\mathcal O_S(M')\) 上忠实平坦，则
\[
  u^t:\mathcal O_S(M')\longrightarrow\mathcal O_S(M)
\]
为单射；当 \(S\ne\varnothing\) 时，这蕴含 \(u:M'\to M\) 为单射。

为证明 \SGARefLink{sga3:viii:corollary:3:2:item:b:01}{(b)}，注意若 \(u\) 为满态射，则 \(\mathcal O_S(M')\to\mathcal O_S(M)\) 为满射，所以 \(u^t\) 为闭浸入，从而更是单态射。反过来，若 \(u^t\) 为单态射，则 \(\operatorname{Ker}u^t\) 为单位群。置 \(M''=\operatorname{Coker}u\)，由前面已知
\[
  \operatorname{Ker}u^t\simeq D_S(M'').
\]
于是由 \SGARefLink{sga3:viii:proposition:2:1:item:d:01}{2.1(d)} 得 \(M''=0\)，所以 \(u\) 为满态射。

通常的论证给出 \SGARefLink{sga3:viii:theorem:3:1}{3.1} 的\SGARefLink{sga3:viii:corollary:3:3}{如下推论}。

\medskip
\noindent\SGARefTarget{sga3:viii:corollary:3:3}\textbf{推论 3.3.}
\textit{设}
\SGARefTarget{sga3:viii:corollary:3:3:diagram:00}
\[
  M'\xrightarrow{u}M\xrightarrow{v}M''
\]
\textit{为通常的交换群的正合列，考虑转置列}
\SGARefTarget{sga3:viii:corollary:3:3:diagram:01}
\[
  G''\xrightarrow{v^t}G\xrightarrow{u^t}G'.
\]
\textit{则 \(v^t\) 诱导从 \(G''\) 到 \(\operatorname{Ker}u^t\) 的忠实平坦拟紧态射，且此核为可对角化群，同构于}
\[
  D_S(v(M))=D_S(\operatorname{Coker}u).
\]

% Source page: Exp8 standalone p. 10; combined-reader p. 626.
\medskip
\noindent\SGARefTarget{sga3:viii:corollary:3:4}\textbf{推论 3.4.}
\textit{设 \(S\) 为预概形，\(u:G\to H\) 为局部可对角化的 \(S\)-群预概形的同态，其中 \(H\) 在 \(S\) 上有限型，并置 \(G'=\operatorname{Ker}u\)。则：}
\begin{enumerate}
  \item[(a)]\SGARefTarget{sga3:viii:corollary:3:4:item:a:01}
  \textit{\(G'\) 局部可对角化，且若 \(G\) 在 \(S\) 上有限型，则它也是；}

  \item[(b)]\SGARefTarget{sga3:viii:corollary:3:4:item:b:01}
  \textit{商 \(G/G'\)“存在”：更确切地说，\(G'\) 在 \(G\) 上定义的等价关系为 \(\mathcal M\)-有效的，其中 \(\mathcal M\) 为忠实平坦拟紧态射族（参见 IV，\SGARefLink{sga3:iv:subsection:3:4}{3.4}）。此外，\(G/G'\) 局部可对角化且在 \(S\) 上有限型；}

  \item[(c)]\SGARefTarget{sga3:viii:corollary:3:4:item:c:01}
  \textit{同态 \(u:G\to H\) 唯一分解为}
  \SGARefTarget{sga3:viii:corollary:3:4:diagram:01}
  \[
    G\xrightarrow{v}G/G'\xrightarrow{w}H,
  \]
  \textit{其中 \(v\) 为典范同态，因而忠实平坦且拟紧。此外，\(w\) 为闭浸入，从而更是单态射。}
\end{enumerate}
\textit{最后，商}
\[
  H'=H/\operatorname{Im}w
    =\operatorname{Coker}w
    =\operatorname{Coker}u
\]
\textit{存在：更确切地说，\(G/G'\) 在 \(H\) 上定义的等价关系为 \(\mathcal M\)-有效的，且 \(H'\) 在 \(S\) 上有限型。}

由商 \(G/G'\) 的定义（参见 IV，\SGARefLink{sga3:iv:lemma:3:2:3}{3.2.3}），\SGARefLink{sga3:viii:corollary:3:4:item:c:01}{(c)} 的第一个断言由 \SGARefLink{sga3:viii:corollary:3:4:item:b:01}{(b)} 得出。
\footnote{编者注：以下扩充了原文。}
来证明 \((\mathrm{fpqc})\) 商层 \(\widetilde G/\widetilde{G'}\) 可表。这与所有其余断言一样，可以在 \(S\) 上局部验证；因此可以假设 \(G\) 和 \(H\) 都可对角化，形如 \(D_S(M)\) 和 \(D_S(N)\)。

由于 \(H\) 在 \(S\) 上有限型，由 \SGARefLink{sga3:viii:proposition:2:1:item:b:01}{2.1(b)}，群 \(N\) 为有限型。因此，由 \SGARefLink{sga3:viii:corollary:1:5}{1.5}，\(u\) 由同态 \(u':N\to M\) 定义。由 \SGARefLink{sga3:viii:theorem:3:1}{3.1} 和 \SGARefLink{sga3:viii:corollary:3:2}{3.2} 得
\[
  G'\simeq D_S(\operatorname{Coker}u')
\]
且 \(\widetilde G/\widetilde{G'}\) 由 \(D_S(\operatorname{Im}u')\) 表示。此外，考虑正合列
\SGARefTarget{sga3:viii:corollary:3:4:diagram:02}
\[
  0\longrightarrow\operatorname{Ker}u'
   \longrightarrow N
   \xrightarrow{w'}\operatorname{Im}u'
   \longrightarrow0,
\]
可知 \(w\) 为闭浸入，且
\[
  H'=H/\operatorname{Im}w=D_S(\operatorname{Ker}u').
\]
后者在 \(S\) 上有限型，因为 \(N\)，从而 \(\operatorname{Ker}u'\)，均为有限型。

\medskip
\noindent\SGARefTarget{sga3:viii:remarks:3:4}\textbf{注。}
\SGARefLink{sga3:viii:corollary:3:4}{3.4} 中商的存在性结果将在 \(\SGARefLink{sga3:viii:section:5}{\S5}\) 中得到大幅推广。

另一方面，注意在\SGARefLink{sga3:viii:section:3}{本节}和\SGARefLink{sga3:viii:section:2}{上一节}中，假设 \(S\ne\varnothing\) 仅用于保证某些逆向结论成立，以便由可对角化 \(S\)-群的某些假设推出相应通常群的性质。“正向”的结果不需要对 \(S\) 作任何限制，这里所给的证明一般均适用。

\medskip
\noindent\SGARefTarget{sga3:viii:corollary:3:5}\textbf{推论 3.5.}
\textit{设 \(G\) 为 \(S\) 上的可对角化群预概形，\(n\ne0\) 为整数。则 \(G\) 中作为同态}
\[
  n\mathbin{\cdot}\operatorname{id}_G:G\longrightarrow G,
\]
\textit{的核的子群 \({}_nG\)，在 \(S\) 上整；若 \(G\) 在 \(S\) 上有限型，它还在 \(S\) 上有限。}

% Source page: Exp8 standalone p. 11; combined-reader p. 627.
事实上，若 \(G=D_S(M)\)，则由 \SGARefLink{sga3:viii:theorem:3:1}{3.1}，
\[
  {}_nG=D_S(M/nM)
\]
结果由 \SGARefLink{sga3:viii:proposition:2:1:item:b:01}{2.1(b)}、\SGARefLink{sga3:viii:proposition:2:1:item:c:01}{(c)} 和 \SGARefLink{sga3:viii:proposition:2:1:item:cprime:01}{\((c')\)} 得出。

\SGARefTarget{sga3:viii:section:4}\section*{4. 可对角化群下的挠子}
\addcontentsline{toc}{section}{4. 可对角化群下的挠子}

\SGARefTarget{sga3:viii:section:4:problem:01}设 \(S\) 为预概形，\(G=D_S(M)\) 为 \(S\) 上的可对角化群。我们将确定 \(S\) 上关于忠实平坦拟紧拓扑的 \(G\)-挠子（torsor，或主齐性 \(G\)-丛）；参见第 IV 讲，\SGARefLink{sga3:iv:subsection:5:1}{5.1}。回顾，若带有 \(S\)-算子群 \(G\) 的 \(S\) 上预概形 \(P\) 满足：\(S\) 的每个点都有开邻域 \(U\) 和忠实平坦拟紧态射 \(S'\to U\)，使得
\[
  P'=P\times_SS'
\]
作为带算子的对象同构于
\[
  G'=G\times_SS',
\]
其中 \(G'\) 通过右平移作用于自身，则称它为\emph{挠子}或\emph{主齐性丛}。由于 \(G\) 在 \(S\) 上仿射，由 SGA 1，VIII 5.6 可知 \(P\) 必在 \(S\) 上仿射。此外，由于 \(G\) 本身在 \(S\) 上忠实平坦且拟紧，\(P\) 在 \(G\) 下为主齐性，当且仅当它是形式主齐性的，且在 \(S\) 上忠实平坦并拟紧（参见 IV，\SGARefLink{sga3:iv:proposition:5:1:6}{5.1.6}）。

还要回顾（\SGARefLink{sga3:i:proposition:4-7-3}{第 I 讲，4.7.3}），给定带有 \(S\)-算子群 \(G=D_S(M)\) 的、在 \(S\) 上仿射的 \(S\)-预概形 \(P\)，相当于给定 \(S\) 上按 \(M\) 分次的交换拟凝聚代数。换言之，相当于给定 \(S\) 上的拟凝聚代数 \(\mathscr A\)，连同模的直和分解
\[
  \mathscr A=\coprod_{m\in M}\mathscr A_m,
\]
满足
\[
  \mathscr A_m\mathbin{\cdot}\mathscr A_{m'}
  \subseteq\mathscr A_{m+m'}
  \qquad\text{对 }m,m'\in M.
\]
在此约定下，\SGARefLink{sga3:viii:section:4:problem:01}{上述问题}的答案如下。

\medskip
\noindent\SGARefTarget{sga3:viii:proposition:4:1}\textbf{命题 4.1.}
\textit{设 \(P\) 为 \(M\)-分次代数 \(\mathscr A\) 所定义的、带有 \(S\)-算子群 \(G=D_S(M)\) 的预概形。则 \(P\) 为 \(G\) 下的主齐性丛，当且仅当 \(\mathscr A\) 满足以下条件：}
\begin{enumerate}
  \item[(a)]\SGARefTarget{sga3:viii:proposition:4:1:item:a:01}
  \textit{对每个 \(m\in M\)，\(\mathscr A_m\) 为 \(S\) 上的可逆模；}

  \item[(b)]\SGARefTarget{sga3:viii:proposition:4:1:item:b:01}
  \textit{对 \(m,m'\in M\)，由 \(\mathscr A\) 的乘法诱导的同态}
  \[
    \mathscr A_m\otimes_{\mathcal O_S}\mathscr A_{m'}
    \longrightarrow\mathscr A_{m+m'}
  \]
  \textit{为同构。}
\end{enumerate}

\SGARefLink{sga3:viii:proposition:4:1}{这些条件}的必要性直接由下降得到，因为当 \(P\) 为平凡主齐性丛，即 \(\mathscr A=\mathcal O_S(M)\) 时，它们成立。至于充分性，注意 \SGARefLink{sga3:viii:proposition:4:1:item:a:01}{(a)} 已蕴含 \(P\) 在 \(S\) 上忠实平坦；它总是在 \(S\) 上拟紧，因为它在 \(S\) 上仿射。因此只需验证 \(P\) 在 \(G\) 下为形式主齐性的，即熟知的态射
\SGARefTarget{sga3:viii:proposition:4:1:diagram:01}
\[
  P\times_SG\longrightarrow P\times_SP
\]
为同构。在仿射代数上，\SGARefLink{sga3:viii:proposition:4:1:diagram:01}{此态射}明确地由同态
\SGARefTarget{sga3:viii:proposition:4:1:diagram:02}
\[
  \mathscr A\otimes\mathscr A
  \longrightarrow
  \mathscr A(M)=\mathscr A\otimes\mathcal O_S(M).
\]
给出。

% Source page: Exp8 standalone p. 12; combined-reader p. 628.
在双次数 \((m,n)\)（其中 \(m,n\in M\)）上，它由
\[
  x_m\otimes y_n\longmapsto x_my_n\otimes e_n.
\]
给出。在次数层面，\SGARefLink{sga3:viii:proposition:4:1:diagram:02}{此同态}与下述同态 \(M\times M\to M\times M\) 相容：
\[
  (m,n)\longmapsto(m+n,n),
\]
后者为同构。这表明，\SGARefLink{sga3:viii:proposition:4:1:item:b:01}{(b)} 恰好表达 \(P\) 为形式主齐性的，且这一点不依赖于 \SGARefLink{sga3:viii:proposition:4:1:item:a:01}{(a)}；这就证明了 \SGARefLink{sga3:viii:proposition:4:1}{4.1}。

忠实平坦下降还给出：

\medskip
\noindent\SGARefTarget{sga3:viii:corollary:4:2}\textbf{推论 4.2.}
\textit{\SGARefLink{sga3:viii:proposition:4:1}{4.1} 的条件蕴含同态}
\[
  \mathcal O_S\longrightarrow\mathscr A_0
\]
\textit{为同构。}

例如，若 \(M=\mathbb Z\)，则在 \SGARefLink{sga3:viii:proposition:4:1}{4.1} 的条件下，代数 \(\mathscr A\) 本质上由 \(\mathscr A_1=\mathscr L\) 确定，即有分次代数的同构
\[
  \mathscr A\simeq\coprod_{n\in\mathbb Z}\mathscr L^{\otimes n}
\]
这样就重新得到熟知的结果：

\medskip
\noindent\SGARefTarget{sga3:viii:corollary:4:3}\textbf{推论 4.3.}
\textit{\(S\) 上 \(\mathbb G_{m,S}\) 下的主齐性丛 \(P\) 的范畴，等价于 \(S\) 上可逆模 \(\mathscr L\) 的范畴，其中两个范畴的态射都是相应结构的同构。互为拟逆的函子如下给出：对 \(P\) 联系其 \(\mathbb Z\)-分次仿射代数的 \(1\) 次分量；对 \(\mathscr L\) 联系下述 \(\mathbb Z\)-分次代数的谱：}
\[
  \coprod_{n\in\mathbb Z}\mathscr L^{\otimes n}.
\]

特别地：

\medskip
\noindent\SGARefTarget{sga3:viii:corollary:4:4}\textbf{推论 4.4.}
\textit{\(S\) 上 \(\mathbb G_{m,S}\) 下的主齐性丛的同构类群，同构于 \(S\) 上可逆模的同构类群 \(\operatorname{Pic}(S)\)，即}
\[
  H^1(S,\mathcal O_S^\times).
\]

由于 \(\mathbb G_{m,S}\) 为模 \(\mathcal O_S\) 的自同构概形，推论 \SGARefLink{sga3:viii:corollary:4:4}{4.4} 等价于\SGARefLink{sga3:viii:corollary:4:5}{下述陈述}；这是 Hilbert“\SGARefLink{sga3:viii:corollary:4:5}{定理 90}”的一种形式。

\medskip
\noindent\SGARefTarget{sga3:viii:corollary:4:5}\textbf{推论 4.5.}
\textit{\(S\) 上 \(\mathbb G_m\) 下的每个主齐性丛，关于 Zariski 拓扑都是局部平凡的。}

\medskip
\noindent\SGARefTarget{sga3:viii:remark:4:5:1}\textbf{注 4.5.1.}
对于 \(\mu_n\) 这样的群，或 \(\mathbb G_m\) 的“扭曲形式”，\SGARefLink{sga3:viii:corollary:4:5}{上述陈述}一般不再成立。例如，实数域 \(\mathbb R\) 上 \(\mathbb G_m\) 的唯一扭曲形式的一阶上同调群等于 \(\mathbb Z/2\mathbb Z\)。

% Source page: Exp8 standalone p. 13; combined-reader p. 629.
\footnote{编者注：增加了以下内容。}
事实上，令 \(\mathbb S^1\) 为范数态射
\[
  N:\prod_{\mathbb C/\mathbb R}\mathbb G_{m,\mathbb C}
  \longrightarrow\mathbb G_{m,\mathbb R};
\]
的核；它是 \(\mathbb G_{m,\mathbb R}\) 的一个 \(\mathbb C/\mathbb R\)-扭曲形式。\(\prod_{\mathbb C/\mathbb R}(\mathbb G_{m,\mathbb C})\) 中的方程 \(N(z)=-1\) 定义了 \(\operatorname{Spec}(\mathbb R)\) 上的 \(\mathbb S^1\)-挠子 \(X\)；它关于 平展 拓扑局部平凡，但不平凡，因为 \(X(\mathbb R)=\varnothing\)。来证明
\[
  H^1_{\mathrm{\acute et}}(\mathbb R,\mathbb S^1)
  \simeq\mathbb Z/2\mathbb Z.
\]
有平滑交换代数 \(\mathbb R\)-群的正合列
\SGARefTarget{sga3:viii:remark:4:5:1:diagram:01}
\[
  1\longrightarrow\mathbb S^1
  \longrightarrow\prod_{\mathbb C/\mathbb R}\mathbb G_{m,\mathbb C}
  \longrightarrow\mathbb G_{m,\mathbb R}
  \longrightarrow1,
\]
它给出 平展 上同调（或 Galois 上同调）的长正合列：
\SGARefTarget{sga3:viii:remark:4:5:1:diagram:02}
\[
\begin{aligned}
  0\longrightarrow\mathbb S^1(\mathbb R)
  &\longrightarrow\mathbb C^\times
   \xrightarrow{N}\mathbb R^\times
   \longrightarrow H^1_{\mathrm{\acute et}}(\mathbb R,\mathbb S^1)\\
  &\longrightarrow
   H^1_{\mathrm{\acute et}}\!
   \left(\mathbb R,
     \prod_{\mathbb C/\mathbb R}\mathbb G_{m,\mathbb C}\right)
   \longrightarrow\cdots .
\end{aligned}
\]
而（例如见 XXIV，\SGARefLink{sga3:xxiv:proposition:8-4}{8.4}），
\[
  H^1_{\mathrm{\acute et}}\!
  \left(\mathbb R,
    \prod_{\mathbb C/\mathbb R}\mathbb G_{m,\mathbb C}\right)
  \simeq
  H^1_{\mathrm{\acute et}}(\mathbb C,\mathbb G_{m,\mathbb C}),
\]
后一个群由 \SGARefLink{sga3:viii:corollary:4:5}{4.5} 为零（在这里也可因为 \(\mathbb C\) 代数闭而得出）。因此得到同构
\[
  H^1_{\mathrm{\acute et}}(\mathbb R,\mathbb S^1)
  \simeq\mathbb R^\times/N(\mathbb C^\times)
  \simeq\{\pm1\}.
\]

\SGARefLink{sga3:viii:section:5}{下一节}将需要\SGARefLink{sga3:viii:proposition:4:6}{以下结果}。

\medskip
\noindent\SGARefTarget{sga3:viii:proposition:4:6}\textbf{命题 4.6.}
\textit{在 \SGARefLink{sga3:viii:proposition:4:1}{4.1} 的条件下，条件 \SGARefLink{sga3:viii:proposition:4:1:item:a:01}{(a)} 和 \SGARefLink{sga3:viii:proposition:4:1:item:b:01}{(b)} 等价于下列条件：}
\begin{enumerate}
  \item[\((a')\)]\SGARefTarget{sga3:viii:proposition:4:6:item:aprime:01}
  \textit{同态 \(\mathcal O_S\to\mathscr A_0\) 为同构；}

  \item[\((b')\)]\SGARefTarget{sga3:viii:proposition:4:6:item:bprime:01}
  \textit{对每个 \(m\in M\)（只要求 \(M\) 的某个生成系中的元素满足就够了），有}
  \[
    \mathscr A_m\mathbin{\cdot}\mathscr A_{-m}=\mathscr A_0.
  \]
\end{enumerate}

由 \SGARefLink{sga3:viii:corollary:4:2}{4.2}，必要性显然。\footnote{编者注：将“\SGARefLink{sga3:viii:proposition:4:1}{4.1}”改成了“\SGARefLink{sga3:viii:corollary:4:2}{4.2}”。}因此归结到证明\SGARefLink{sga3:viii:corollary:4:7}{下面的推论}。

\medskip
\noindent\SGARefTarget{sga3:viii:corollary:4:7}\textbf{推论 4.7.}
\textit{设}
\[
  A=\coprod_{n\in\mathbb Z}A_n
\]
\textit{为 \(\mathbb Z\)-分次环，满足}
\[
  A_1\mathbin{\cdot}A_{-1}=A_0.
\]
\textit{则各 \(A_n\) 为可逆 \(A_0\)-模，且对 \(n,n'\in\mathbb Z\)，由 \(A\) 的乘法诱导的同态}
\[
  A_n\otimes_{A_0}A_{n'}\longrightarrow A_{n+n'}
\]
\textit{为同构。}

由假设，存在 \(f_i\in A_1\) 和 \(g_i\in A_{-1}\)，使得
\SGARefTarget{sga3:viii:corollary:4:7:eq:star:01}
\[
  \makebox[\displaywidth][s]{%
    \(\text{(*)}\)\hfill
    \(\displaystyle\sum_i f_ig_i=1.\)\hfill
    \(\phantom{\text{(*)}}\)%
  }
\]
由于待证断言关于 \(\operatorname{Spec}(A_0)\) 是局部的，\footnote{编者注：扩充了以下一段。}且 \(\SGARefLink{sga3:viii:corollary:4:7:eq:star:01}{(*)}\) 表明 \(\operatorname{Spec}(A_0)\) 被仿射开子集 \(D(f_ig_i)\) 覆盖，故归结到存在元素 \(f\in A_1\) 在 \(A\) 中可逆的情形。于是对每个 \(n\in\mathbb Z\)，\(f^n\) 为 \(A_n\) 中的可逆元，因而定义同构
\[
  A_0\xrightarrow{\ \sim\ }A_n,
  \qquad h\longmapsto f^nh.
\]
此外，
